\chapter{The Last Apple Tree}

The cow climbed steadily, hauling the cart through a narrow rift between rocky ridges. Exposed roots thrust out from ancient stumps, bleached and brittle. They followed a dry creekbed that twisted upward, the stones grinding beneath the wheels.

Mist gathered as they ascended. At the ridge crest the cow paused, breathing heavily, looking back. A few drops of poisoned rain stirred the leaves. Wind passed through the dead trees along the ridge, whispering through their skeletal branches. Tibor urged the cow onward.

Beyond the ridge lay a rocky field overgrown with plantain and dandelion, choked with the desiccated stalks of former weeds. A ruined fence sagged ahead. Tibor consulted his map, holding it up like a sacred text. Yes—this was correct. He would encounter the southern tribes here.

The cart broke through the fence and halted near a collapsed well. Nearby lay the remains of a structure: sagging beams, broken glass, rusted springs, scraps of cloth. At the field’s edge stood an orchard of dead trees—row upon row of lifeless trunks, stripped and blackened, bent by endless wind.

One tree remained alive.

It was an apple tree, gnarled and barren, bearing only a few dark leaves and several withered apples clinging stubbornly to the branches. Tibor could not look away. The sight repelled and fascinated him. All the others had failed. This one endured.

He plucked a leaf and examined it. The branches creaked together. The sound unsettled him.

Night approached rapidly. Cold wind struck him, and the valley below dissolved into shadow. Leaves tore loose and spun past. One nearly brushed his face. Exhaustion swept over him. Fear followed.

He turned to flee.

Then he saw the apple.

It glistened faintly in the dark, red beneath the blackened mist. He told himself it was rotten, worthless—yet it seemed to beckon. He activated his radio and called for Father Handy. Only static answered.

Lost, hungry, alone, Tibor felt his resolve falter. Perhaps this was a magic tree, like those Father Handy spoke of. Christians would call the apple evil, a test of sin. But the Servants of Wrath rejected that myth. And he had eaten nothing all day.

He decided to pick it up—but not eat it.

As he lifted the apple into the light, movement flickered at the edge of his vision.

Two figures approached.

They were tall—nearly eight feet—thin, ash-blue, their skin horny and scaled. Mutants: lizard-type. One carried an ax. The other wore only tattered cloth. Their eyes were large, curious.

They greeted him politely.

They introduced themselves as Jackson and Earl Potter. They examined his cart with fascination, shook hands with his metal grippers, commented on the cow. Tibor masked his aversion and spoke carefully. He claimed to be from Charlottesville, exaggerating its size.

They spoke of hunting, of brass blow-dart weapons that followed prey unerringly once launched. They invited him to their settlement. He declined, explaining his mission.

He showed them the photograph of Carleton Lufteufel.

They did not recognize it.

A third lizard joined them, studying Tibor intently. It quickly deduced that Tibor was on a pilgrimage—and then revealed something disturbing: another human was following him. A young man on a bicycle.

Pete Sands.

Tibor’s fear sharpened. He asked for concealment. The lizards offered help—and then something else.

A dog.

They produced the animal: large, gray, affectionate. It whined joyfully at Tibor. The sight stirred something deep and painful in him. He wanted it desperately.

The lizards insisted. Humans must be protected, they said. For the future. For intact genes.

Before Tibor could answer, the lizards stiffened.

Bugs had arrived.

Five of them emerged into the light—upright, chitin-shelled, buzzing voices sharp with contempt. They insulted the lizards crudely, boasting and jeering. The lizards relaxed, however, amused rather than threatened.

Then Tibor saw why.

Behind the bugs stood a score of runners.

The balance of the night shifted.

